\documentclass[aspectratio=169]{beamer}

\usetheme{Madrid}
\usecolortheme{default}

\usepackage[spanish,es-noshorthands]{babel}
\usepackage[utf8]{inputenc}
\usepackage[T1]{fontenc}
\usepackage{lmodern}
\usepackage{amsmath,amssymb}
\usepackage{array}
\usepackage{xcolor}
\usepackage{booktabs}

\definecolor{USTGreen}{RGB}{27,94,32}
\definecolor{USTLight}{RGB}{232,245,233}
\definecolor{USTDark}{RGB}{18,68,25}

\setbeamercolor{structure}{fg=USTGreen}
\setbeamercolor{frametitle}{fg=white,bg=USTGreen}
\setbeamercolor{title}{fg=white,bg=USTGreen}
\setbeamercolor{block title}{fg=white,bg=USTGreen}
\setbeamercolor{block body}{bg=USTLight,fg=black}
\setbeamertemplate{navigation symbols}{}
\setbeamertemplate{blocks}[rounded][shadow=false]

\title[Sistemas de ecuaciones lineales]{Sistemas de Ecuaciones Lineales con Dos Incógnitas}
\subtitle{Clasificación, métodos algebraicos e interpretación en contextos agronómicos}
\author{Fundamentos de Matemática Básica}
\institute{Universidad Santo Tomás}
\date{}

\begin{document}
	
%------------------------------------------------
\begin{frame}
	\titlepage
\end{frame}

%------------------------------------------------
\begin{frame}{Objetivo general}
	Analizar y resolver sistemas de ecuaciones lineales con dos incógnitas mediante métodos algebraicos y geométricos, interpretando el significado de sus soluciones en situaciones propias de la agronomía.
	
	\bigskip
	
	\begin{block}{Resultados de aprendizaje}
		\begin{itemize}
			\item Identificar sistemas compatibles determinados, compatibles indeterminados e incompatibles.
			\item Resolver sistemas mediante sustitución, igualación y reducción.
			\item Interpretar geométricamente y técnicamente la solución de un sistema.
			\item Detectar patrones, relaciones e inconsistencias en datos agrícolas.
		\end{itemize}
	\end{block}
\end{frame}

%------------------------------------------------
\begin{frame}{Motivación agronómica}
	En agronomía es frecuente trabajar con dos variables que se relacionan simultáneamente.
	
	\bigskip
	
	\begin{block}{Ejemplos de variables}
		\begin{itemize}
			\item Superficie sembrada con dos cultivos distintos.
			\item Cantidad de dos fertilizantes en una mezcla.
			\item Litros de agua aplicados por riego y por fertirriego.
			\item Jornadas de maquinaria y jornadas de mano de obra.
		\end{itemize}
	\end{block}
	
	\bigskip
	
	Un sistema permite modelar restricciones simultáneas, por ejemplo: presupuesto, superficie, dosis, volumen o productividad.
\end{frame}
	
%------------------------------------------------
\begin{frame}{Definición}
	Un sistema de ecuaciones lineales con dos incógnitas tiene la forma:
	
	\[
	\begin{cases}
		ax+by=c\\
		dx+ey=f
	\end{cases}
	\]
	
	donde:
	
	\[
	a,b,c,d,e,f\in\mathbb{R}
	\]
	
	\bigskip
	
	Resolver el sistema significa encontrar los valores de \(x\) e \(y\) que satisfacen simultáneamente ambas ecuaciones.
	
	\bigskip
	
	\begin{block}{Lectura en contexto}
		Si \(x\) representa hectáreas de maíz e \(y\) hectáreas de trigo, la solución debe tener sentido agronómico: no puede haber hectáreas negativas y las unidades deben ser coherentes.
	\end{block}
\end{frame}
	
%------------------------------------------------
\begin{frame}{Interpretación geométrica}
	Cada ecuación lineal representa una recta en el plano cartesiano.
	
	\bigskip
	
	Por lo tanto, resolver un sistema equivale a estudiar la relación entre dos rectas.
	
	\bigskip
	
	\begin{center}
		\begin{tabular}{|c|c|c|}
			\hline
			\textbf{Caso} & \textbf{Rectas} & \textbf{Soluciones} \\
			\hline
			Compatible determinado & Se cortan & Una solución \\
			\hline
			Compatible indeterminado & Coinciden & Infinitas soluciones \\
			\hline
			Incompatible & Paralelas & Sin solución \\
			\hline
		\end{tabular}
	\end{center}
\end{frame}
	
%------------------------------------------------
\begin{frame}{Sistema compatible determinado}
	Tiene una única solución.
	
	\[
	\begin{cases}
		x+y=5\\
		2x-y=4
	\end{cases}
	\]
	
	\smallskip
	
	Sumamos ambas ecuaciones:
	
\[
\begin{array}{rcl}
	(x+y)+(2x-y) &=& 5+4 \\[4pt]
	3x &=& 9 \\[4pt]
	x &=& 3
\end{array}
\]
	
	Luego:
	
	\[
	\begin{array}{rcl}
	3+y &=& 5 \\
	y &=& 2
	\end{array}
	\]
	
	\[
	\boxed{(x,y)=(3,2)}
	\]
\end{frame}

%------------------------------------------------
\begin{frame}{Ejemplo agronómico: sistema compatible determinado}
	\small
	Un predio destina \(5\) hectáreas a dos cultivos: maíz y trigo. Sea:
	
	\[
	x=\text{hectáreas de maíz}, \qquad y=\text{hectáreas de trigo}
	\]
	
	La superficie total es \(5\) ha. Además, por una restricción de fertilización nitrogenada, se obtiene el modelo:
	
	\[
	\begin{cases}
		x+y=5\\
		2x-y=4
	\end{cases}
	\]
	
	\begin{block}{Interpretación}
		La solución \((3,2)\) indica que se deben sembrar \(3\) ha de maíz y \(2\) ha de trigo. Es un sistema compatible determinado porque existe una única planificación que satisface ambas restricciones.
	\end{block}
\end{frame}
	
%------------------------------------------------
\begin{frame}{Sistema compatible indeterminado}
	Tiene infinitas soluciones.
	
	\[
	\begin{cases}
		x+y=4\\
		2x+2y=8
	\end{cases}
	\]
	
	La segunda ecuación es múltiplo de la primera:
	
	\[
	2(x+y)=8
	\]
	
	\[
	x+y=4
	\]
	
	\bigskip
	
	Ambas ecuaciones representan la misma recta.
	
	\[
	\boxed{\text{Infinitas soluciones}}
	\]
\end{frame}

%------------------------------------------------
\begin{frame}{Ejemplo agronómico: sistema compatible indeterminado}
	\small
	Un técnico registra la cantidad de dos enmiendas aplicadas en una parcela:
	
	\[
	x=\text{toneladas de compost}, \qquad y=\text{toneladas de guano}
	\]
	
	Dos informes indican:
	
	\[
	\begin{cases}
		x+y=4\\
		2x+2y=8
	\end{cases}
	\]
	
	\begin{block}{Análisis}
		Los informes entregan exactamente la misma información, porque la segunda ecuación es solo el doble de la primera. Por lo tanto, cualquier combinación que sume \(4\) toneladas cumple ambas condiciones.
	\end{block}
	
	\[
	\boxed{\text{Sistema compatible indeterminado}}
	\]
\end{frame}
	
%------------------------------------------------
\begin{frame}{Sistema incompatible}
	No tiene solución.
	
	\[
	\begin{cases}
		x+y=3\\
		x+y=7
	\end{cases}
	\]
	
	\bigskip
	
	La misma expresión \(x+y\) no puede ser igual a 3 y a 7 al mismo tiempo.
	
	\bigskip
	
	Esto genera una contradicción:
	
	\[
	3=7
	\]
	
	\[
	\boxed{\text{Sistema incompatible}}
	\]
\end{frame}

%------------------------------------------------
\begin{frame}{Ejemplo agronómico: sistema incompatible}
	\small
	En una planificación de riego, \(x\) representa metros cúbicos aplicados por goteo e \(y\) metros cúbicos aplicados por aspersión.
	
	\bigskip
	
	Dos registros del mismo día indican:
	
	\[
	\begin{cases}
		x+y=3\\
		x+y=7
	\end{cases}
	\]
	
	\begin{block}{Interpretación técnica}
		La cantidad total de agua aplicada no puede ser simultáneamente \(3\) m\(^3\) y \(7\) m\(^3\). Esto evidencia una inconsistencia en los datos o en el registro del predio.
	\end{block}
	
	\[
	\boxed{\text{No existe solución válida}}
	\]
\end{frame}
	
%------------------------------------------------
\begin{frame}{Método de sustitución}
	Consiste en despejar una incógnita en una ecuación y reemplazarla en la otra.
	
	\bigskip
	
	Resolver:
	
	\[
	\begin{cases}
		x+y=10\\
		x=y+2
	\end{cases}
	\]
	
	Sustituimos \(x=y+2\):
	
	\[
	(y+2)+y=10
	\]
	
	\[
	2y+2=10
	\]
	
	\[
	2y=8
	\]
	
	\[
	y=4
	\]
	
	\[
	x=4+2=6
	\]
	
	\[
	\boxed{(x,y)=(6,4)}
	\]
\end{frame}

%------------------------------------------------
\begin{frame}{Ejemplo agronómico por sustitución}
	\small
	Para preparar \(120\) litros de solución nutritiva se mezclan agua y concentrado. Sea:
	
	\[
	x=\text{litros de agua}, \qquad y=\text{litros de concentrado}
	\]
	
	La receta indica que el agua debe ser cinco veces el concentrado:
	
	\[
	\begin{cases}
		x+y=120\\
		x=5y
	\end{cases}
	\]
	
	\begin{align*}
	5y+y&=120\\
	6y&=120\\
	y&=20\\
	x&=5(20)=100
	\end{align*}
	
	\[
	\boxed{x=100\text{ L de agua},\qquad y=20\text{ L de concentrado}}
	\]
\end{frame}
	
%------------------------------------------------
\begin{frame}{Método de igualación}
	Consiste en despejar la misma incógnita en ambas ecuaciones e igualar las expresiones.
	
	\bigskip
	
	Resolver:
	
	\[
	\begin{cases}
		x=8-y\\
		x=2y+2
	\end{cases}
	\]
	
	Igualamos:
	
	\[
	8-y=2y+2
	\]
	
	\[
	6=3y
	\]
	
	\[
	y=2
	\]
	
	Luego:
	
	\[
	x=8-2=6
	\]
	
	\[
	\boxed{(x,y)=(6,2)}
	\]
\end{frame}

%------------------------------------------------
\begin{frame}{Ejemplo agronómico por igualación}
	\small
	En un invernadero se distribuyen bandejas de lechuga entre dos sectores. Sea:
	
	\[
	x=\text{bandejas en el sector A}, \qquad y=\text{bandejas en el sector B}
	\]
	
	Dos criterios operativos entregan:
	
	\[
	\begin{cases}
		x=90-2y\\
		x=y+15
	\end{cases}
	\]
	
	\begin{align*}
	90-2y&=y+15\\
	75&=3y\\
	y&=25
	\end{align*}
	
	\[
	x=y+15=25+15=40
	\]
	
	\[
	\boxed{x=40\text{ bandejas},\qquad y=25\text{ bandejas}}
	\]
\end{frame}
	
%------------------------------------------------
\begin{frame}{Método de reducción}
	Consiste en eliminar una incógnita sumando o restando ecuaciones.
	
	\bigskip
	
	Resolver:
	
	\[
	\begin{cases}
		2x+y=11\\
		x-y=1
	\end{cases}
	\]
	
	Sumamos:
	
	\[
	(2x+y)+(x-y)=11+1
	\]
	
	\[
	3x=12
	\]
	
	\[
	x=4
	\]
	
	Sustituimos:
	
	\[
	4-y=1
	\]
	
	\[
	y=3
	\]
	
	\[
	\boxed{(x,y)=(4,3)}
	\]
\end{frame}

%------------------------------------------------
\begin{frame}{Ejemplo agronómico por reducción}
	\small
	Se desea mezclar dos fertilizantes líquidos A y B. Sea:
	
	\[
	x=\text{litros del fertilizante A}, \qquad y=\text{litros del fertilizante B}
	\]
	
	El aporte total de nitrógeno y fósforo queda modelado por:
	
	\[
	\begin{cases}
		30x+20y=260\\
		10x+20y=140
	\end{cases}
	\]
	
	Restamos la segunda ecuación a la primera:
	
	\begin{align*}
	(30x+20y)-(10x+20y)&=260-140\\
	20x&=120\\
	x&=6
	\end{align*}
	
	Sustituyendo:
	
	\[
	10(6)+20y=140\quad \Rightarrow\quad 20y=80\quad \Rightarrow\quad y=4
	\]
	
	\[
	\boxed{x=6\text{ L de A},\qquad y=4\text{ L de B}}
	\]
\end{frame}
	
%------------------------------------------------
\begin{frame}{Ejercicio 1: fertirriego por sustitución}
	\small
	Un estanque debe preparar \(72\) litros de mezcla entre agua y fertilizante soluble. La cantidad de agua debe ser el triple de la cantidad de fertilizante.
	
	\bigskip
	
	Sea:
	
	\[
	x=\text{litros de agua}, \qquad y=\text{litros de fertilizante soluble}
	\]
	
	\[
	\begin{cases}
		x+y=72\\
		x=3y
	\end{cases}
	\]
	
	\begin{block}{Preguntas}
		\begin{itemize}
			\item Resuelva el sistema por sustitución.
			\item Interprete la solución en litros.
			\item Verifique si las cantidades cumplen la proporción indicada.
		\end{itemize}
	\end{block}
\end{frame}
	
%------------------------------------------------
\begin{frame}{Solución ejercicio 1}
	\small
	Sistema:
	
	\[
	\begin{cases}
		x+y=72\\
		x=3y
	\end{cases}
	\]
	
	Sustituimos \(x=3y\):
	
	\begin{align*}
	3y+y&=72\\
	4y&=72\\
	y&=18
	\end{align*}
	
	\[
	x=3(18)=54
	\]
	
	\[
	\boxed{x=54\text{ L de agua},\qquad y=18\text{ L de fertilizante}}
	\]
	
	\begin{block}{Verificación}
		\(54+18=72\) y \(54=3\cdot18\). El sistema es compatible determinado.
	\end{block}
\end{frame}
	
%------------------------------------------------
\begin{frame}{Ejercicio 2: bandejas de almácigos por igualación}
	\small
	En un vivero, la cantidad de bandejas ubicadas en el sector A puede describirse de dos formas según restricciones de espacio y manejo:
	
	\[
	x=\text{bandejas en el sector A}, \qquad y=\text{bandejas en el sector B}
	\]
	
	\[
	\begin{cases}
		x=90-2y\\
		x=y+15
	\end{cases}
	\]
	
	\begin{block}{Tarea}
		Resolver por igualación, clasificar el sistema e interpretar la solución en el contexto del vivero.
	\end{block}
\end{frame}
	
%------------------------------------------------
\begin{frame}{Solución ejercicio 2}
	\small
	Igualamos las expresiones para \(x\):
	
	\begin{align*}
	90-2y&=y+15\\
	75&=3y\\
	y&=25
	\end{align*}
	
	Luego:
	
	\[
	x=y+15=25+15=40
	\]
	
	\[
	\boxed{(x,y)=(40,25)}
	\]
	
	\begin{block}{Interpretación}
		Se deben ubicar \(40\) bandejas en el sector A y \(25\) bandejas en el sector B. El sistema es compatible determinado.
	\end{block}
\end{frame}
	
%------------------------------------------------
\begin{frame}{Ejercicio 3: costos agrícolas por reducción}
	\small
	Un predio compra dos tipos de insumos: bioestimulante y corrector de suelo. Sea:
	
	\[
	x=\text{unidades de bioestimulante}, \qquad y=\text{unidades de corrector de suelo}
	\]
	
	Dos presupuestos parciales quedan representados por:
	
	\[
	\begin{cases}
		4x+2y=540\\
		3x+2y=450
	\end{cases}
	\]
	
	\begin{block}{Tarea}
		Resolver por reducción, verificar la solución e interpretar las cantidades obtenidas.
	\end{block}
\end{frame}
	
%------------------------------------------------
\begin{frame}{Solución ejercicio 3}
	\small
	Restamos la segunda ecuación a la primera:
	
	\begin{align*}
	(4x+2y)-(3x+2y)&=540-450\\
	x&=90
	\end{align*}
	
	Sustituimos en la segunda ecuación:
	
	\begin{align*}
	3(90)+2y&=450\\
	270+2y&=450\\
	2y&=180\\
	y&=90
	\end{align*}
	
	\[
	\boxed{(x,y)=(90,90)}
	\]
	
	\begin{block}{Verificación}
		\(4(90)+2(90)=540\) y \(3(90)+2(90)=450\). El sistema es compatible determinado.
	\end{block}
\end{frame}
	
%------------------------------------------------
\begin{frame}{Patrones e inconsistencias}
	Durante la resolución pueden aparecer patrones importantes:
	
	\bigskip
	
	\begin{block}{Contradicción}
		Si aparece una igualdad falsa:
		
		\[
		0=5
		\]
		
		el sistema es incompatible.
	\end{block}
	
	\begin{block}{Identidad}
		Si aparece una igualdad verdadera:
		
		\[
		0=0
		\]
		
		el sistema es compatible indeterminado.
	\end{block}
\end{frame}

%------------------------------------------------
\begin{frame}{Patrones en datos agrícolas}
	\small
	En agronomía, estos patrones pueden revelar información relevante:
	
	\begin{block}{Caso 1: identidad}
		Dos restricciones de fertilización son equivalentes. Esto indica que falta una condición adicional para determinar una única mezcla.
	\end{block}
	
	\begin{block}{Caso 2: contradicción}
		Dos registros de riego, superficie o dosis se contradicen. Esto indica error de medición, transcripción o interpretación técnica.
	\end{block}
	
	\begin{block}{Caso 3: solución única}
		Las restricciones son independientes y permiten tomar una decisión concreta: cuánto sembrar, cuánto aplicar o cuánto comprar.
	\end{block}
\end{frame}
	
%------------------------------------------------
\begin{frame}{Actividad de clasificación contextualizada}
	\small
	Clasifique cada sistema e interprete su significado agrícola.
	
	\[
	A)
	\begin{cases}
		x+y=10\\
		2x+2y=20
	\end{cases}
	\]
	
	\[
	B)
	\begin{cases}
		x+y=10\\
		x+y=15
	\end{cases}
	\]
	
	\[
	C)
	\begin{cases}
		x+y=10\\
		x-y=2
	\end{cases}
	\]
	
	\begin{block}{Contexto común}
		\(x\): hectáreas destinadas a cultivo A. \quad \(y\): hectáreas destinadas a cultivo B.
	\end{block}
\end{frame}
	
%------------------------------------------------
\begin{frame}{Solución actividad de clasificación}
	\small
	\[
	A)
	\begin{cases}
		x+y=10\\
		2x+2y=20
	\end{cases}
	\]
	
	Compatible indeterminado: ambas ecuaciones expresan la misma superficie total.
	
	\bigskip
	
	\[
	B)
	\begin{cases}
		x+y=10\\
		x+y=15
	\end{cases}
	\]
	
	Incompatible: el mismo predio no puede tener simultáneamente \(10\) ha y \(15\) ha asignadas.
	
	\bigskip
	
	\[
	C)
	\begin{cases}
		x+y=10\\
		x-y=2
	\end{cases}
	\]
	
	Compatible determinado: existe una única distribución de hectáreas.
\end{frame}

%------------------------------------------------
\begin{frame}{Ejercicios propuestos: agronomía aplicada}
	\small
	\begin{enumerate}
		\item \textbf{Superficie de cultivos.} Un predio tiene \(18\) ha entre hortalizas y frutales. La superficie de frutales es el doble de la superficie de hortalizas.
		\[
		\begin{cases}
		x+y=18\\
		x=2y
		\end{cases}
		\]
		Resolver por sustitución.
		
		\item \textbf{Planificación de riego.} Dos criterios técnicos para el volumen del sector A son:
		\[
		\begin{cases}
		x=20-2y\\
		x=y+5
		\end{cases}
		\]
		Resolver por igualación.
	\end{enumerate}
\end{frame}

%------------------------------------------------
\begin{frame}{Ejercicios propuestos: reducción y clasificación}
	\small
	\begin{enumerate}
		\setcounter{enumi}{2}
		\item \textbf{Compra de insumos.} Un agricultor compra dos productos A y B. Las restricciones de costo se modelan por:
		\[
		\begin{cases}
		3x+2y=24\\
		x-2y=0
		\end{cases}
		\]
		Resolver por reducción.
		
		\item \textbf{Control de consistencia.} Clasifique los sistemas e indique si representan datos coherentes o inconsistentes:
		\[
		\begin{cases}
		x+y=8\\
		3x+3y=24
		\end{cases}
		\qquad
		\begin{cases}
		2x-y=5\\
		4x-2y=12
		\end{cases}
		\]
	\end{enumerate}
\end{frame}
	
%------------------------------------------------
\begin{frame}{Cierre}
	\begin{block}{Ideas clave}
		\begin{itemize}
			\item Un sistema lineal con dos incógnitas representa dos rectas y dos restricciones simultáneas.
			\item Puede tener una solución, infinitas soluciones o ninguna solución.
			\item Los métodos algebraicos principales son sustitución, igualación y reducción.
			\item La solución debe verificarse algebraicamente y analizarse en contexto agronómico.
			\item En aplicaciones reales, una contradicción puede indicar datos mal registrados o restricciones imposibles.
		\end{itemize}
	\end{block}
\end{frame}
	
%------------------------------------------------
\begin{frame}{Tarea contextualizada}
	\small
	Resolver y clasificar los siguientes sistemas. En cada caso, indique la interpretación agronómica de la solución.
	
	\[
	1)
	\begin{cases}
		2x+y=15\\
		x-y=3
	\end{cases}
	\qquad
	\begin{array}{l}
	 x=\text{hectáreas de cultivo A}\vspace{1mm}\\
	 y=\text{hectáreas de cultivo B}
	\end{array}
	\]
	
	\[
	2)
	\begin{cases}
		x+y=8\\
		3x+3y=24
	\end{cases}
	\qquad
	\begin{array}{l}
	 x=\text{toneladas de compost}\vspace{1mm}\\
	 y=\text{toneladas de guano}
	\end{array}
	\]
	
	\[
	3)
	\begin{cases}
		2x-y=5\\
		4x-2y=12
	\end{cases}
	\qquad
	\begin{array}{l}
	 x=\text{m}^{3}\text{ de agua de pozo}\vspace{1mm}\\
	 y=\text{m}^{3}\text{ de agua acumulada}
	\end{array}
	\]
\end{frame}
	
\end{document}
